% Gerado por ferramentas/gerar.py; não edite à mão.
\section*{Exercícios}
\addcontentsline{toc}{section}{Exercícios}
Os pontos antes de cada número indicam a dificuldade: \nivel{1} fácil, \nivel{2} médio, \nivel{3} difícil e \nivel{4} desafio. Os exercícios marcados com \estrela{} têm solução completa, com animação do gráfico, no repositório do livro: \url{https://github.com/AleffAzevedo/calculo-1-pt/blob/main/exercicios/cap4/README.md}. As respostas estão no fim do capítulo.

\begin{itemize}[leftmargin=5.2em,labelsep=0.6em,itemsep=4pt]
\item[\exnum{1}{4.1}] Encontre a linearização de $f(x)=\sqrt[3]x$ em $a=8$ e use-a para estimar $\sqrt[3]{8{,}2}$.
\item[\exnum{1}{4.2}] Para $y=x^2-3x$, compare $dy$ e $\Delta y$ quando $x=2$ e $dx=0{,}1$.
\item[\exnum{1}{4.3}] Encontre os pontos críticos de $f(x)=x^3-12x+1$.
\item[\exnum{1}{4.4}] Encontre os pontos críticos de $f(x)=x^{2/3}(x-5)$ (inclua os pontos em que $f'$ não existe).
\item[\exnum{1}{4.5}] Encontre os extremos absolutos de $f(x)=x^2-4x+3$ em $[0,3]$.
\item[\exnum{1}{4.6}] Encontre os extremos absolutos de $f(x)=\sin x+\cos x$ em $[0,\pi]$.
\item[\exnum{1}{4.7}] Verifique as hipóteses do Teorema de Rolle para $f(x)=x^2-4x$ em $[0,4]$ e encontre $c$.
\item[\exnum{1}{4.8}] Encontre o $c$ do TVM para $f(x)=x^2$ em $[1,3]$ e interprete.
\item[\exnum{1}{4.9}] Determine os intervalos de crescimento e de decrescimento de $f(x)=x^3-3x^2$.
\item[\exnum{1}{4.10}] Estude a concavidade de $f(x)=x^4-4x^3$ e encontre os pontos de inflexão.
\item[\exnum{1}{4.11}] Use o teste da segunda derivada para classificar os pontos críticos de $f(x)=2x^3-3x^2-12x$.
\item[\exnum{1}{4.12}] Calcule $\displaystyle\lim_{x\to\infty}\frac{3x^2-x}{5x^2+2}$.
\item[\exnum{1}{4.13}] Calcule $\displaystyle\lim_{x\to-\infty}\frac{2x+1}{\sqrt{x^2+3}}$.
\item[\exnum{1}{4.14}] Encontre as assíntotas verticais e horizontais de $f(x)=\dfrac{2x^2}{x^2-9}$.
\item[\exnum{1}{4.15}] Encontre a assíntota oblíqua de $f(x)=\dfrac{x^2+1}{x+1}$.
\item[\exnum{1}{4.16}] Use L'Hôpital para calcular $\displaystyle\lim_{x\to1}\frac{\ln x}{x-1}$ e $\displaystyle\lim_{x\to0}\frac{e^x-1}{x}$.
\item[\exnum{1}{4.17}] Calcule $\displaystyle\lim_{x\to\infty}x^2e^{-x}$.
\item[\exnum{1}{4.18}] Faça duas iterações do método de Newton para $f(x)=x^2-5$ com $x_0=2$.
\item[\exnum{1}{4.19}] Calcule $\int(3x^2-4x+5)\,dx$ e $\int(\sin x+e^x)\,dx$.
\item[\exnum{1}{4.20}] Resolva o problema de valor inicial $f'(x)=6x^2-2$, $f(1)=3$.
\item[\exnum{2}{4.21}] \estrela\ \textbf{A escada escorregando.} Uma escada de 5 m está apoiada numa parede vertical. A base se afasta da parede a
$0{,}5$ m/s. (a) Com que rapidez o topo desce quando a base está a 3 m da parede?
(b) Em que posição as duas extremidades se movem com a mesma rapidez?
\item[\exnum{2}{4.22}] \estrela\ \textbf{A sombra no poste.} Uma pessoa de $1{,}8$ m se afasta de um poste de luz de 6 m de altura a $1{,}5$ m/s.
(a) Com que rapidez sua sombra aumenta? (b) Com que rapidez a ponta da sombra se move?
\item[\exnum{2}{4.23}] \estrela\ \textbf{Avião e radar.} Um avião voa horizontalmente a 4 km de altitude, a 600 km/h, numa trajetória que passa
diretamente acima de uma estação de radar. Com que rapidez a distância avião–radar aumenta
quando essa distância é 5 km (e o avião já passou da estação)?
\item[\exnum{2}{4.24}] Uma mancha de óleo circular se espalha com o raio crescendo a 2 m/s. Com que rapidez a área cresce quando $r=10$ m?
\item[\exnum{2}{4.25}] A aresta de um cubo cresce a $0{,}5$ cm/s. Com que rapidez o volume cresce quando a aresta mede 4 cm?
\item[\exnum{2}{4.26}] O raio de uma esfera mede $10$ cm com erro de $\pm0{,}2$ cm. Estime o erro absoluto e o percentual no volume.
\item[\exnum{2}{4.27}] Use uma aproximação linear para estimar $\sin31^\circ$.
\item[\exnum{2}{4.28}] \estrela\ \textbf{Método do intervalo fechado.} Encontre os valores máximo e mínimo absolutos de $f(x)=x\,e^{-x}$ em $[0,4]$.
\item[\exnum{2}{4.29}] \estrela\ \textbf{Teorema do Valor Médio.} Verifique as hipóteses do TVM para $f(x)=x^3-x$ em $[0,2]$ e encontre todos os $c$ garantidos
pelo teorema.
\item[\exnum{2}{4.30}] Mostre que $x^3+x-1=0$ tem exatamente uma raiz real.
\item[\exnum{2}{4.31}] Mostre que $\tan^{-1}x+\tan^{-1}\tfrac1x=\tfrac\pi2$ para $x>0$.
\item[\exnum{2}{4.32}] Encontre os extremos locais de $f(x)=x^2e^{-x}$.
\item[\exnum{2}{4.33}] Encontre os extremos locais de $f(x)=x-2\sin x$ em $(0,2\pi)$.
\item[\exnum{2}{4.34}] Calcule $\displaystyle\lim_{x\to\infty}\big(\sqrt{x^2+x}-x\big)$.
\item[\exnum{2}{4.35}] Calcule $\displaystyle\lim_{x\to0}\frac{x-\sin x}{x^3}$.
\item[\exnum{2}{4.36}] Calcule $\displaystyle\lim_{x\to\infty}x^{1/x}$.
\item[\exnum{2}{4.37}] Calcule $\displaystyle\lim_{x\to\pi/2^-}(\sec x-\tan x)$.
\item[\exnum{2}{4.38}] \estrela\ \textbf{O limite $(1+a/x)^x$.} Mostre, com L'Hôpital, que $\displaystyle\lim_{x\to\infty}\Big(1+\frac ax\Big)^x=e^a$ para
qualquer constante $a$. Use o resultado para calcular o montante de R\$\,1000 a 6\% ao ano
com capitalização contínua por 10 anos.
\item[\exnum{2}{4.39}] Calcule $\displaystyle\int\frac{x^3-2x}{x^2}\,dx$.
\item[\exnum{2}{4.40}] Calcule $\displaystyle\int\Big(\frac2{1+x^2}-\frac3{\sqrt{1-x^2}}\Big)dx$.
\item[\exnum{3}{4.41}] \estrela\ \textbf{Esboço completo de um gráfico.} Faça o estudo completo e esboce $f(x)=\dfrac{x}{x^2+1}$: domínio, simetria, assíntotas,
crescimento, extremos, concavidade e inflexões.
\item[\exnum{3}{4.42}] \estrela\ \textbf{Cercado junto ao rio.} Um fazendeiro tem 600 m de cerca para construir um cercado retangular junto a um rio reto
(sem cerca do lado do rio), dividido em dois currais por uma cerca perpendicular ao rio. Que
dimensões maximizam a área total?
\item[\exnum{3}{4.43}] \estrela\ \textbf{Ponto mais próximo.} Encontre os pontos da parábola $y=x^2$ mais próximos do ponto $(0,3)$ e a distância mínima.
\item[\exnum{3}{4.44}] \estrela\ \textbf{O cabo mais barato.} Uma usina fica numa margem de um rio de 1 km de largura. Uma fábrica fica na outra
margem, 3 km rio abaixo. O cabo custa R\$\,5 mil/km sob a água e R\$\,4 mil/km em terra.
Qual é o trajeto mais barato (atravessar em diagonal até um ponto e seguir pela margem)?
\item[\exnum{3}{4.45}] \estrela\ \textbf{Método de Newton para $\cos x=x$.} Use o método de Newton com $x_0=1$ para aproximar a solução de $\cos x=x$ com 9 casas
decimais. Quantas iterações são necessárias?
\item[\exnum{3}{4.46}] \estrela\ \textbf{Lançamento vertical.} Uma bola é lançada para cima a partir de 2 m de altura com velocidade de 15 m/s
($g=9{,}8$ m/s²). Partindo de $a(t)=-9{,}8$, encontre $v(t)$ e $s(t)$, a altura máxima e o
instante em que a bola atinge o chão.
\item[\exnum{3}{4.47}] Uma calha de 10 m tem seção em triângulo isósceles invertido, com 1 m de largura no topo e $0{,}5$ m de altura. Entra água a $0{,}4$ m³/min. Com que rapidez o nível sobe quando $h=0{,}2$ m?
\item[\exnum{3}{4.48}] Dois carros partem do mesmo ponto: um para o norte a 60 km/h e outro para o leste a 80 km/h. Com que rapidez a distância entre eles aumenta?
\item[\exnum{3}{4.49}] Um balão sobe verticalmente a 2 m/s e é observado de um ponto no solo a 50 m da base. Com que rapidez o ângulo de elevação aumenta quando o balão está a 50 m de altura?
\item[\exnum{3}{4.50}] Encontre o retângulo de maior área inscrito num semicírculo de raio $R$ (com um lado sobre o diâmetro).
\item[\exnum{3}{4.51}] Uma caixa sem tampa de base quadrada deve ter volume $32\,000$ cm³. Que dimensões minimizam o material?
\item[\exnum{3}{4.52}] Encontre o retângulo de maior área com base no eixo $x$ e vértices superiores na parábola $y=12-x^2$.
\item[\exnum{3}{4.53}] A demanda de um produto é $q=1000-20p$ ($p$ em reais). Que preço maximiza a receita $R=pq$?
\item[\exnum{3}{4.54}] Faça o estudo completo de $f(x)=x^{2/3}(x-5)$ (há uma cúspide) e esboce o gráfico.
\item[\exnum{3}{4.55}] Faça o estudo completo de $f(x)=e^{-x^2}$ e esboce o gráfico.
\item[\exnum{3}{4.56}] Calcule $\displaystyle\lim_{x\to0}\frac{e^x-1-x}{x^2}$ e $\displaystyle\lim_{x\to\infty}x\Big(\frac\pi2-\tan^{-1}x\Big)$.
\item[\exnum{3}{4.57}] Use o TVM para provar que $|\sin a-\sin b|\le|a-b|$ para todos $a,b$.
\item[\exnum{3}{4.58}] Prove que $e^x\ge1+x$ para todo $x$.
\item[\exnum{3}{4.59}] Resolva $f''(x)=6x+2$, $f'(0)=1$, $f(0)=-1$.
\item[\exnum{3}{4.60}] Prove que a equação $\cos x=x$ tem exatamente uma solução.
\item[\exnum{4}{4.61}] Prove que a única função derivável com $f'=f$ e $f(0)=1$ é $f(x)=e^x$. (Dica: derive $g(x)=f(x)e^{-x}$.)
\item[\exnum{4}{4.62}] A luz vai de $A$ (no ar, velocidade $v_1$) a $B$ (na água, velocidade $v_2$) pelo caminho de tempo mínimo. Deduza a lei de Snell $\dfrac{\sin\theta_1}{v_1}=\dfrac{\sin\theta_2}{v_2}$.
\item[\exnum{4}{4.63}] Para $a>0$, mostre que o método de Newton para $x^2=a$ dá $x_{n+1}=\tfrac12\big(x_n+\tfrac a{x_n}\big)$ e que $x_n\ge\sqrt a$ para todo $n\ge1$, qualquer que seja $x_0>0$.
\item[\exnum{4}{4.64}] Calcule $\displaystyle\lim_{x\to0}\Big(\frac1{x^2}-\frac1{\sin^2x}\Big)$.
\item[\exnum{4}{4.65}] Qual é o comprimento da maior barra que pode ser carregada horizontalmente pela esquina entre dois corredores de larguras 8 m e 27 m?
\end{itemize}
\subsection*{Respostas}
{\small\begin{multicols}{2}\raggedright
\begin{itemize}[leftmargin=2.7em,labelsep=0.5em,itemsep=1pt]
\item[\textbf{4.1}] $L(x)=2+\tfrac{x-8}{12}$; $\approx2{,}0167$.
\item[\textbf{4.2}] $dy=0{,}1$; $\Delta y=0{,}11$.
\item[\textbf{4.3}] $x=\pm2$.
\item[\textbf{4.4}] $x=2$ e $x=0$.
\item[\textbf{4.5}] Máximo $3$ em $x=0$; mínimo $-1$ em $x=2$.
\item[\textbf{4.6}] Máximo $\sqrt2$ em $x=\tfrac\pi4$; mínimo $-1$ em $x=\pi$.
\item[\textbf{4.7}] $c=2$.
\item[\textbf{4.8}] $c=2$: tangente paralela à secante.
\item[\textbf{4.9}] Cresce em $(-\infty,0)\cup(2,\infty)$; decresce em $(0,2)$.
\item[\textbf{4.10}] Para cima em $(-\infty,0)\cup(2,\infty)$, para baixo em $(0,2)$; inflexões $(0,0)$ e $(2,-16)$.
\item[\textbf{4.11}] Máx. local $f(-1)=7$; mín. local $f(2)=-20$.
\item[\textbf{4.12}] $\tfrac35$.
\item[\textbf{4.13}] $-2$.
\item[\textbf{4.14}] Verticais $x=\pm3$; horizontal $y=2$.
\item[\textbf{4.15}] $y=x-1$.
\item[\textbf{4.16}] $1$ e $1$.
\item[\textbf{4.17}] $0$.
\item[\textbf{4.18}] $x_1=2{,}25$; $x_2\approx2{,}236111$.
\item[\textbf{4.19}] $x^3-2x^2+5x+C$; $-\cos x+e^x+C$.
\item[\textbf{4.20}] $f(x)=2x^3-2x+3$.
\item[\textbf{4.21}] (a) $0{,}375$ m/s; (b) quando $x=y=\tfrac{5}{\sqrt2}\approx3{,}54$ m.
\item[\textbf{4.22}] (a) $\tfrac{9}{14}\approx0{,}64$ m/s; (b) $\tfrac{15}{7}\approx2{,}14$ m/s.
\item[\textbf{4.23}] 360 km/h.
\item[\textbf{4.24}] $40\pi\approx125{,}7$ m²/s.
\item[\textbf{4.25}] $24$ cm³/s.
\item[\textbf{4.26}] $dV=80\pi\approx251$ cm³; $6\%$.
\item[\textbf{4.27}] $\tfrac12+\tfrac{\sqrt3}2\cdot\tfrac\pi{180}\approx0{,}5151$.
\item[\textbf{4.28}] Máximo $\tfrac1e\approx0{,}368$ em $x=1$; mínimo $0$ em $x=0$.
\item[\textbf{4.29}] $c=\tfrac{2}{\sqrt3}\approx1{,}155$.
\item[\textbf{4.30}] Existe pelo TVI; é única porque $f'=3x^2+1>0$ (pelo Rolle, duas raízes exigiriam $f'=0$).
\item[\textbf{4.31}] A derivada da soma é $0$, então ela é constante; em $x=1$ vale $\tfrac\pi2$.
\item[\textbf{4.32}] Mín. local $f(0)=0$; máx. local $f(2)=\tfrac4{e^2}\approx0{,}541$.
\item[\textbf{4.33}] Mín. $f(\tfrac\pi3)=\tfrac\pi3-\sqrt3$; máx. $f(\tfrac{5\pi}3)=\tfrac{5\pi}3+\sqrt3$.
\item[\textbf{4.34}] $\tfrac12$.
\item[\textbf{4.35}] $\tfrac16$.
\item[\textbf{4.36}] $1$.
\item[\textbf{4.37}] $0$.
\item[\textbf{4.38}] $e^a$; $1000e^{0{,}6}\approx$ R\$\,1822{,}12.
\item[\textbf{4.39}] $\tfrac{x^2}2-2\ln|x|+C$.
\item[\textbf{4.40}] $2\tan^{-1}x-3\sin^{-1}x+C$.
\item[\textbf{4.41}] Ímpar; assíntota $y=0$; máx. $(1,\tfrac12)$, mín. $(-1,-\tfrac12)$; inflexões em $x=0$ e $x=\pm\sqrt3$.
\item[\textbf{4.42}] 100 m (perpendicular ao rio) por 300 m (paralelo); área 30\,000 m².
\item[\textbf{4.43}] $\big(\pm\sqrt{5/2},\,\tfrac52\big)$; distância $\tfrac{\sqrt{11}}{2}\approx1{,}658$.
\item[\textbf{4.44}] Atravessar até o ponto a $\tfrac43$ km rio abaixo; custo mínimo R\$\,15 mil.
\item[\textbf{4.45}] $x\approx0{,}739085133$; 3 iterações (a 4ª só confirma).
\item[\textbf{4.46}] $s(t)=-4{,}9t^2+15t+2$; altura máx. $\approx13{,}48$ m em $t\approx1{,}53$ s; chão em $t\approx3{,}19$ s.
\item[\textbf{4.47}] $0{,}1$ m/min.
\item[\textbf{4.48}] $100$ km/h (constante).
\item[\textbf{4.49}] $0{,}02$ rad/s.
\item[\textbf{4.50}] Lados $\sqrt2R$ e $\tfrac R{\sqrt2}$; área $R^2$.
\item[\textbf{4.51}] Base $40\times40$ cm, altura $20$ cm; área $4800$ cm².
\item[\textbf{4.52}] Base 4, altura 8; área 32.
\item[\textbf{4.53}] $p=25$; receita R\$\,12\,500.
\item[\textbf{4.54}] Máx. local $(0,0)$ (cúspide); mín. local $(2,-3\sqrt[3]4)\approx(2;-4{,}76)$; inflexão $(-1,-6)$.
\item[\textbf{4.55}] Par; máx. $(0,1)$; inflexões em $x=\pm\tfrac1{\sqrt2}$; assíntota $y=0$.
\item[\textbf{4.56}] $\tfrac12$ e $1$.
\item[\textbf{4.57}] $\sin a-\sin b=\cos c\,(a-b)$ e $|\cos c|\le1$.
\item[\textbf{4.58}] $g(x)=e^x-1-x$ tem mínimo absoluto $g(0)=0$.
\item[\textbf{4.59}] $f(x)=x^3+x^2+x-1$.
\item[\textbf{4.60}] Existência pelo TVI; unicidade porque $g(x)=\cos x-x$ tem $g'=-\sin x-1\le0$, zerando só em pontos isolados.
\item[\textbf{4.61}] $g'=0$, logo $g$ é constante $=g(0)=1$.
\item[\textbf{4.62}] Escreva o tempo em função do ponto de cruzamento e anule a derivada.
\item[\textbf{4.63}] Desigualdade entre as médias aritmética e geométrica.
\item[\textbf{4.64}] $-\tfrac13$.
\item[\textbf{4.65}] $13\sqrt{13}\approx46{,}9$ m (minimizando $\tfrac8{\sin\theta}+\tfrac{27}{\cos\theta}$).
\end{itemize}\end{multicols}}
